% ===========================================================================
\section{分解纲领：单群与合成列}\label{sec:structure}

\subsection{群的分解纲领：从阿贝尔群到非阿贝尔群}

\block{回顾（阿贝尔群）}
阿贝尔群可以分解为若干循环群的积；每个循环群又可分解为若干
\keyword{素数幂阶循环群}的积。所以任何有限阿贝尔群都可以分解成素数幂阶
循环群的积——对阿贝尔群而言，结构问题基本就清楚了。

\block{非阿贝尔群的纲领}
对非阿贝尔群 $G$，一般的目标是得到一列正规子群
\[
  \{e\} \normal G_1 \normal G_2 \normal G_3 \normal G,
\]
有了这列正规子群，就可以把问题分解成一串“因子”：
\[
  G_1,\quad G_2/G_1,\quad G_3/G_2,\quad G/G_3,
\]
即把群分解成若干\keyword{商群}分别研究。我们希望把它分解成“最简”的形式——
这就引出单群的概念。

\subsection{单群与合成列}

\begin{definition}[单群（Simple Group）]
群 $G$ 称为\keyword{单群}，如果它没有非平凡的正规子群，即
\[
  N \normal G \;\Longrightarrow\; N = G \ \text{或}\ N = \{e\}.
\]
\end{definition}

\begin{example}
素数阶循环群都是单群。
\end{example}

\begin{theorem}[合成列的存在性]
任何\keyword{有限群} $G$ 都存在一列如上的正规子群，使得每个因子
$G_{i+1}/G_i$ 都是单群。
\end{theorem}

\begin{proof}
假如某个因子不是单群，比如 $G/G_3$ 不是单群，那么它有一个非平凡正规子群
$\bar N$（既不等于 $G/G_3$，也不是平凡群）。由\keyword{商群的对应定理}，
$\bar N$ 在自然同态 $\pi: G\to G/G_3$ 下的\keyword{原像}
$\pi^{-1}(\bar N)$ 是 $G$ 的正规子群，且严格夹在 $G_3$ 与 $G$ 之间。
把这个中间群“切进”子群列中，得到更细的分解；有限群保证这一过程在有限步
后终止，归纳即得。
\end{proof}
